% Options for packages loaded elsewhere
\PassOptionsToPackage{unicode}{hyperref}
\PassOptionsToPackage{hyphens}{url}
%
\documentclass[
]{article}
\usepackage{amsmath,amssymb}
\usepackage{iftex}
\ifPDFTeX
  \usepackage[T1]{fontenc}
  \usepackage[utf8]{inputenc}
  \usepackage{textcomp} % provide euro and other symbols
\else % if luatex or xetex
  \usepackage{unicode-math} % this also loads fontspec
  \defaultfontfeatures{Scale=MatchLowercase}
  \defaultfontfeatures[\rmfamily]{Ligatures=TeX,Scale=1}
\fi
\usepackage{lmodern}
\ifPDFTeX\else
  % xetex/luatex font selection
\fi
% Use upquote if available, for straight quotes in verbatim environments
\IfFileExists{upquote.sty}{\usepackage{upquote}}{}
\IfFileExists{microtype.sty}{% use microtype if available
  \usepackage[]{microtype}
  \UseMicrotypeSet[protrusion]{basicmath} % disable protrusion for tt fonts
}{}
\makeatletter
\@ifundefined{KOMAClassName}{% if non-KOMA class
  \IfFileExists{parskip.sty}{%
    \usepackage{parskip}
  }{% else
    \setlength{\parindent}{0pt}
    \setlength{\parskip}{6pt plus 2pt minus 1pt}}
}{% if KOMA class
  \KOMAoptions{parskip=half}}
\makeatother
\usepackage{xcolor}
\usepackage[margin=1in]{geometry}
\usepackage{graphicx}
\makeatletter
\def\maxwidth{\ifdim\Gin@nat@width>\linewidth\linewidth\else\Gin@nat@width\fi}
\def\maxheight{\ifdim\Gin@nat@height>\textheight\textheight\else\Gin@nat@height\fi}
\makeatother
% Scale images if necessary, so that they will not overflow the page
% margins by default, and it is still possible to overwrite the defaults
% using explicit options in \includegraphics[width, height, ...]{}
\setkeys{Gin}{width=\maxwidth,height=\maxheight,keepaspectratio}
% Set default figure placement to htbp
\makeatletter
\def\fps@figure{htbp}
\makeatother
\setlength{\emergencystretch}{3em} % prevent overfull lines
\providecommand{\tightlist}{%
  \setlength{\itemsep}{0pt}\setlength{\parskip}{0pt}}
\setcounter{secnumdepth}{-\maxdimen} % remove section numbering
\usepackage[spanish]{babel}
\usepackage{amsmath}
\ifLuaTeX
  \usepackage{selnolig}  % disable illegal ligatures
\fi
\usepackage{bookmark}
\IfFileExists{xurl.sty}{\usepackage{xurl}}{} % add URL line breaks if available
\urlstyle{same}
\hypersetup{
  hidelinks,
  pdfcreator={LaTeX via pandoc}}

\author{}
\date{\vspace{-2.5em}}

\begin{document}

\includegraphics{01.png}

\includegraphics{02.png}

\includegraphics{03.png}

\section{Pregunta}\label{pregunta}

Camilo desea saber cuánta agua se necesita para llenar el cilindro
interno, pero solamente cuenta con las medidas de las dimensiones que
muestra la figura.

Para calcular cuánta agua se necesita para llenar el cilindro interno,
necesitamos conocer el volumen del espacio interior del cilindro, el
cual se calcula con la fórmula del volumen de un cilindro
\(V = \pi r^2 h\), donde \(r\) es el radio y \(h\) la altura del
cilindro.

\begin{enumerate}
\def\labelenumi{\arabic{enumi}.}
\item
  \textbf{Determinación del radio interno}: En el problema, se indica
  que el radio interno es de 1 metro.
\item
  \textbf{Determinación de la altura del cilindro}: No se especifica
  directamente la altura del cilindro, pero asumiremos que es conocida
  para el cálculo.
\item
  \textbf{Cálculo del volumen}: Usando el radio interno de 1 metro, y
  una altura \(h\) (que asumimos conocida), el volumen del cilindro se
  calcula como: \[
  V = \pi \times (1 \, \text{m})^2 \times h = \pi \times h \, \text{m}^3
  \]
\end{enumerate}

Este volumen representa la cantidad de agua necesaria para llenar el
cilindro hasta su altura \(h\).\\
\strut \\

\subsubsection{\texorpdfstring{\textbf{Solución
Completa}}{Solución Completa}}\label{soluciuxf3n-completa}

\paragraph{\texorpdfstring{\textbf{Paso 1: Identificar las dimensiones
conocidas del
cilindro}}{Paso 1: Identificar las dimensiones conocidas del cilindro}}\label{paso-1-identificar-las-dimensiones-conocidas-del-cilindro}

Conocemos el radio interno del cilindro, que es de 1 metro.

\paragraph{\texorpdfstring{\textbf{Paso 2: Aplicar la fórmula del
volumen de un
cilindro}}{Paso 2: Aplicar la fórmula del volumen de un cilindro}}\label{paso-2-aplicar-la-fuxf3rmula-del-volumen-de-un-cilindro}

Utilizamos la fórmula del volumen \(V = \pi r^2 h\), sustituyendo
\(r = 1 \text{ m}\) y \(h\) (conocida).

\paragraph{\texorpdfstring{\textbf{Respuesta
Final}}{Respuesta Final}}\label{respuesta-final}

El volumen de agua necesaria para llenar el cilindro es
\(\pi h \text{ m}^3\), donde \(h\) es la altura del cilindro.

\subsubsection{\texorpdfstring{\textbf{Concepto
Clave}}{Concepto Clave}}\label{concepto-clave}

\textbf{Volumen del cilindro}: El volumen \(V\) de un cilindro se
calcula como el producto del área de la base circular (que es
\(\pi r^2\)) por la altura \(h\) del cilindro.

\subsubsection{\texorpdfstring{\textbf{Conceptos y ejemplos
relacionados}}{Conceptos y ejemplos relacionados}}\label{conceptos-y-ejemplos-relacionados}

\begin{enumerate}
\def\labelenumi{\arabic{enumi}.}
\tightlist
\item
  \textbf{Área de la base de un cilindro}: Área = \(\pi r^2\).
\item
  \textbf{Perímetro de la base de un cilindro}: Perímetro = \(2\pi r\).
\item
  \textbf{Ejemplo relacionado}: Si un cilindro tiene un radio de 3
  metros y una altura de 5 metros, ¿cuál es su volumen?
  \(V = \pi \times 3^2 \times 5 = 45\pi \text{ m}^3\).
\end{enumerate}

\subsubsection{\texorpdfstring{\textbf{¿Cuál es la opción
correcta?}}{¿Cuál es la opción correcta?}}\label{cuuxe1l-es-la-opciuxf3n-correcta}

Para determinar cuál es la opción correcta, necesitamos analizar cuál de
las opciones dadas es la más relevante para calcular el volumen del
cilindro interno con la información disponible en el problema. Las
opciones son:

A. Radio externo. B. Diámetro externo. C. Altura del cilindro. D.
Perímetro del cilindro.

El cálculo del volumen de un cilindro se realiza mediante la fórmula
\(V = \pi r^2 h\), donde \(r\) es el radio y \(h\) es la altura. Para el
cilindro interno, el radio interno es conocido (1 metro), pero
necesitamos la altura del cilindro para completar el cálculo del
volumen.

\textbf{La opción correcta es:} C. Altura del cilindro.

Esta es la medida crucial faltante que necesitamos para calcular el
volumen de agua que el cilindro puede contener, dado que el radio
interno ya es conocido.

\subsection{\texorpdfstring{\textbf{Retroalimentación Saber
ICFES}}{Retroalimentación Saber ICFES}}\label{retroalimentaciuxf3n-saber-icfes}

\subsubsection{Análisis del Problema}\label{anuxe1lisis-del-problema}

Para resolver este problema, debemos determinar qué medidas son
necesarias para calcular la cantidad de agua que se necesita para llenar
el cilindro interno.

El volumen del cilindro interno se puede calcular utilizando la fórmula
del volumen de un cilindro:

\[ V = \pi r^2 h \]

donde \(r\) es el radio del cilindro y \(h\) es la altura del cilindro.

\subsubsection{Información
Proporcionada}\label{informaciuxf3n-proporcionada}

De la figura y la descripción:

\begin{itemize}
\tightlist
\item
  El radio interno del cilindro es \(1 \,  \text{m}\).
\item
  El grosor del cilindro es \(0.5 \,   \text{m}\).
\item
  No se proporciona directamente la altura del cilindro.
\end{itemize}

\subsubsection{Determinación de las
Opciones}\label{determinaciuxf3n-de-las-opciones}

Vamos a analizar cada opción para determinar cuál es la necesaria para
calcular el volumen del cilindro interno.

\begin{enumerate}
\def\labelenumi{\arabic{enumi}.}
\tightlist
\item
  \textbf{Opción A: Radio externo.}

  \begin{itemize}
  \tightlist
  \item
    El radio externo se puede calcular sumando el grosor al radio
    interno: \(r_{externo} = r_{interno} +  \text{grosor}\).
  \item
    En este caso,
    \(r_{externo} = 1 \,    \text{m} + 0.5 \,   \text{m} = 1.5 \,   \text{m}\).
  \item
    \textbf{No es necesaria esta medida para calcular el volumen del
    cilindro interno.}
  \end{itemize}
\item
  \textbf{Opción B: Diámetro externo.}

  \begin{itemize}
  \tightlist
  \item
    El diámetro externo se puede calcular multiplicando el radio externo
    por 2: \(d_{externo} = 2     \times r_{externo}\).
  \item
    \textbf{No es necesaria esta medida para calcular el volumen del
    cilindro interno.}
  \end{itemize}
\item
  \textbf{Opción C: Altura del cilindro.}

  \begin{itemize}
  \tightlist
  \item
    La altura del cilindro es un componente crucial en la fórmula del
    volumen: \(V = \pi r^2 h\).
  \item
    \textbf{Esta medida es necesaria para calcular el volumen del
    cilindro interno.}
  \end{itemize}
\item
  \textbf{Opción D: Perímetro del cilindro.}

  \begin{itemize}
  \tightlist
  \item
    El perímetro de la base del cilindro (circunferencia) se calcula con
    la fórmula \(P = 2 \pi r\).
  \item
    \textbf{No es necesaria esta medida para calcular el volumen del
    cilindro interno.}
  \end{itemize}
\end{enumerate}

\subsubsection{Retroalimentación}\label{retroalimentaciuxf3n}

\textbf{Nivel de Desempeño:} Nivel 2 - \textbf{Justificación:} El
problema requiere la identificación de las medidas necesarias para
calcular el volumen del cilindro, lo cual implica comprensión básica de
fórmulas geométricas.

\textbf{Competencia:} Formulación y ejecución - \textbf{Justificación:}
El problema implica identificar las dimensiones correctas para aplicar
una fórmula matemática.

\textbf{Componente:} Geometría - \textbf{Justificación:} Se relaciona
con el cálculo de volúmenes de figuras geométricas.

\textbf{Tipo de Pensamiento:} Espacial-métrico - \textbf{Justificación:}
Implica entender las dimensiones y relaciones métricas del cilindro.

\textbf{Estándar Asociado y Categoría:} - \textbf{Estándar Asociado:}
Utiliza propiedades y fórmulas geométricas para resolver problemas. -
\textbf{Categoría:} Contenido no genérico - \textbf{Justificación:}
Aplica conocimientos específicos de geometría.

\textbf{Situaciones o Contexto:} Contexto cotidiano -
\textbf{Justificación:} Se enmarca en una situación práctica de uso de
un cilindro.

\textbf{Problema Analizado:} - \textbf{Habilidad o Conocimiento
Evaluado:} Cálculo de volúmenes utilizando fórmulas geométricas. -
\textbf{Justificación:} Evalúa la capacidad de identificar las
dimensiones necesarias para aplicar una fórmula geométrica.

\textbf{Eje Axial Disciplinar:} Geometría - \textbf{Justificación:} Se
centra en la comprensión y manipulación de figuras geométricas
tridimensionales.

\subsubsection{¿Qué Evalúa?}\label{quuxe9-evaluxfaa}

El problema evalúa la habilidad para identificar las dimensiones
necesarias y aplicar fórmulas geométricas para calcular volúmenes.

\subsubsection{Respuesta Correcta}\label{respuesta-correcta}

\textbf{La respuesta correcta es:} C. Altura del cilindro.

\subsubsection{Justificación de la Respuesta
Correcta}\label{justificaciuxf3n-de-la-respuesta-correcta}

Para calcular el volumen del cilindro interno, necesitamos la altura del
cilindro. El radio interno ya se proporciona, pero la altura es la
medida crucial faltante.

\subsubsection{Motivos que Llevan a Escoger Cada Una de las Opciones
Incorrectas}\label{motivos-que-llevan-a-escoger-cada-una-de-las-opciones-incorrectas}

\begin{itemize}
\tightlist
\item
  \textbf{Opción A (Radio externo):} No es necesario para calcular el
  volumen interno.
\item
  \textbf{Opción B (Diámetro externo):} No es necesario para el volumen
  del cilindro interno.
\item
  \textbf{Opción D (Perímetro del cilindro):} No es útil para calcular
  el volumen del cilindro interno.
\end{itemize}

\subsubsection{Desde Qué Nivel (Grado) de Enseñanza se Puede Aplicar
Este
Ejercicio}\label{desde-quuxe9-nivel-grado-de-enseuxf1anza-se-puede-aplicar-este-ejercicio}

\textbf{Nivel Educativo:} Grado 9 o 10 - \textbf{Justificación:} A este
nivel, los estudiantes deberían tener conocimientos suficientes de
geometría y volúmenes de cuerpos sólidos.

\begin{center}\rule{0.5\linewidth}{0.5pt}\end{center}

Muy bien, aquí tengo un problema. Camilo quiere saber cuánta agua se
necesita para llenar el cilindro interno, pero sólo tiene algunas
medidas de la figura proporcionada. Veamos lo que tiene y lo que le
falta. Primero, el problema menciona ``Diámetro externo'' que es el
diámetro externo, y se da como\ldots{} espera, en realidad, en el
contenido, dice ``Diámetro externo'' pero no da un valor para él
directamente. En su lugar, da el ``Grosor'', que es de 0,5 metros, y el
``Radio interno'', que es de 1 metro. También menciona ``Altura'' y
``Radio externo'', pero no proporciona sus valores. Entonces, Camilo
quiere averiguar el volumen de agua necesario para llenar el cilindro
interno. Para hallar el volumen de un cilindro, la fórmula es V = πr²h,
donde r es el radio y h es la altura. En este caso, como se trata de un
cilindro interior, tiene el radio interior, que es 1 metro. Pero también
necesita la altura del cilindro para calcular el volumen. Mirando las
opciones proporcionadas: 1. Radio externo (radio exterior) Radio externo
2. Diámetro externo (external diameter) 3. Altura del cilindro 4.
Perímetro del cilindro Camilo ya tiene el radio interno y tiene el
espesor, que es de 0,5 metros. Si conoce el radio interno y el grosor,
podría hallar el radio externo sumando el grosor al radio interno
(R\_externo = R\_interno + grosor = 1m + 0,5m = 1,5m). Del mismo modo,
podría hallar el diámetro exterior duplicando el radio exterior
(D\_exterior = 2 * R\_exterior = 3 m). El perímetro del cilindro, o la
circunferencia, no se necesita directamente para calcular el volumen,
pero está relacionado con el radio o diámetro mediante C = 2πr. Lo único
que no tiene es la altura del cilindro. Sin la altura, no puede calcular
el volumen, aunque conozca el radio o el diámetro. Así que, entre las
opciones, la que falta y es necesaria para calcular el volumen es la
``Altura del cilindro''. Permíteme que vuelva a comprobarlo para
asegurarme de que no se me escapa nada. Si tuviera el diámetro exterior,
podría encontrar el radio exterior, pero ya tiene el radio interior y el
grosor, así que es redundante. El perímetro no es necesario para
calcular el volumen. Por lo tanto, sí, la altura es lo que falta.

\end{document}
